\section{The two-state structural alternative}
A certificate that a supplied stochastic machine has the correct
response is not a certificate that a smaller machine cannot have that
response, or that its approximation error is optimal. This distinction
is particularly important for a finite experiment with externally
chosen commands. Each internal state must support all possible future
commands, and a random mixture of whole machines is a larger machine
unless its selector is retained in the charged state.

The two-state component develops an exact global error alternative in the
smallest nontrivial register class. For antisymmetric binary targets,
two labels at every cut reduce without loss to a bounded decomposable
tensor. The seed, each command epoch, and the terminal query are its
modes. The scalar reduction and its parity obstruction are retained
from the preceding revision; their proofs appear in
Section~\ref{sec:binary47}. The present question is how to determine
the distance when all signs are consistent but magnitudes are not.

\begin{theorem}[Sign--magnitude certification]\label{thm:main49}
For every finite antisymmetric binary-response table, the optimal
two-label uniform error has the following complete description.
At each proposed mean tolerance, threshold-active entries determine
all factor coordinates which must be nonzero. A finite sign system
and one bounded logarithmic linear system per sign chamber decide
feasibility. Each rejected chamber has either an impossible positive
endpoint or a nonnegative integer multiplicative certificate on at
most $v+1$ inequalities, where $v=|V_\delta|$ is the number of essential factor coordinates at the tested threshold $\delta$. The sparsity bound is per rejected chamber.
For real-algebraic table entries, univariate root isolation applied
to these certificates determines the optimum and an algebraic
stochastic machine attaining it.
\end{theorem}

Theorem~\ref{thm:main49} is proved in
Theorems~\ref{thm:duality49} and~\ref{thm:algebraic49}.
Its quantifiers are global over the prescribed two-state class.
It does not require a supplied factorization, a strictly positive
target, or a positive optimum. It is not a claim that the union of
chambers is convex or that searching them is polynomial in the word
horizon. The lower certificates are particularly transparent:
weighted multiplication of interval constraints cancels every
unknown factor and produces a violated inequality between data.

The exact example in Theorem~\ref{thm:cubic49} has two different
commands, one command epoch, two positive seeds and their negatives,
and two queries. Its positive-seed table has rational means
$1/10$, $2/5$, and $4/5$. Its optimal binary-TV error is $t/2$, where
\[
 500t^3-375t^2+540t-124=0,\qquad 13/50<t<261/1000.
\]
A five-entry certificate proves the lower bound and explicit
stochastic rows prove the upper bound. The root has algebraic degree
three, so rational tables do not imply a rational optimal machine.
This concerns attainment of a positive error at fixed two-state
width, not the rationality of an exact nonnegative factorization.
The orthogonal example in Proposition~\ref{prop:orthogonal49} gives
an additional arbitrary-horizon exact error with rational witnesses.
The two-bottleneck compatibility frontier from the preceding paper
is included separately in Theorem~\ref{thm:frontier47}.

Fixed-sign matrix approximation and monomial logarithms are classical
\cite{GS19,BKVH07}; global Chebyshev approximation also has a direct
optimization literature \cite{MSZ23}. We use these mechanisms rather
than claiming to originate them. Their precise relation to the
multi-epoch bounded tensor, the sign boundary, and the stochastic
compiler is discussed in Section~\ref{sec:encoding49}. The claimed
finite alternatives are accompanied by proofs, rational certificate
verification and worked exact optima, not by a numerical assertion
that a local search has found the global minimum.

The article keeps a single operational object throughout: one reset,
a finite external command word, and one selected terminal binary
query. The complete prior positive-Hankel, all-word Gram, moment,
grid, arithmetic-width, and finite-bit results follow the new
sections. In particular, the retained entropy and arithmetic growth
claims are not reproved by the finite two-state theorem. The new
rational-height lemma specifies the bit cost of candidate Gram
verification, while the global alternative deliberately states its
larger explicit-table input. Free epoch, horizon, transition tables
and exact real row sampling are distinguished from charged persistent
labels in Definition~\ref{def:machine}.

The orthogonal-word target is equation~\eqref{eq:target}, stated in the introduction.
The new two-state alternative also allows arbitrary finite odd response
 tables; such tables are not silently assumed to have an orthogonal
 realization of a fixed dimension.
